% ============================================================
%  part5/ch27-clio-reconstruction.tex
%  Chapter 27: CLIO Reconstruction
% ============================================================

\chapter{CLIO: Constraint-Leveraged Inference and Optimisation}
\label{ch:clio}

\begin{WhyChapter}
  Part~I established that observation is projection and
  reconstruction is imperfect.
  Part~II derived the reconstruction error $E(x) = d(x,G(F(x)))$
  and showed it is zero only when $F$ is injective.
  This chapter formalises the machinery for minimising
  reconstruction error under admissibility constraints.
  CLIO is the name for that minimisation framework.
\end{WhyChapter}

% Input the core reconstruction propositions from Part II
% ============================================================
%  part2/proofs-part2.tex
%  Core propositions for Part II — Reconstruction Theory
%  These are referenced throughout Chapters 8-11 and Part V
% ============================================================

% This file is \input{} from ch27-clio-reconstruction.tex
% so the propositions appear in context.
% It is kept separate for maintenance.

\section{The reconstruction propositions}
\label{sec:reconstruction-propositions}

We collect here the fundamental results about projection,
reconstruction, and error that form the mathematical core of
the CLIO framework.

\begin{proposition}[Perfect reconstruction criterion]
  \label{prop:perfect-recon}
  $E(x) = 0$ for all $x$ if and only if $F$ is injective
  and $G \circ F = \mathrm{id}$.
\end{proposition}

\begin{proof}
  $(\Rightarrow)$ Suppose $E(x) = d(x, G(F(x))) = 0$ for all $x$.
  Since $d$ is a metric, $G(F(x)) = x$ for all $x$,
  so $G \circ F = \mathrm{id}$.
  If $F(x_1) = F(x_2)$, applying $G$ gives $x_1 = x_2$,
  so $F$ is injective.

  $(\Leftarrow)$ If $F$ is injective and $G \circ F = \mathrm{id}$,
  then $E(x) = d(x, G(F(x))) = d(x,x) = 0$.
\end{proof}

\begin{proposition}[Kernel lower bound]
  \label{prop:kernel-lower}
  Let $[x] = F^{-1}(F(x))$.
  If $|[x]| > 1$, then $E(x) > 0$ for every deterministic
  reconstruction map $G$.
\end{proposition}

\begin{proof}
  If $|[x]| > 1$, there exists $y \neq x$ with $F(y) = F(x)$.
  A deterministic $G$ must satisfy $G(F(x)) = G(F(y))$.
  Hence $G$ cannot simultaneously return $x$ and $y$.
  At least one incurs positive reconstruction error.
\end{proof}

\begin{theorem}[Reconstruction lower bound]
  \label{thm:recon-lower}
  If $F : \mathcal{M} \to L$ is not injective, then no
  reconstruction map $G : L \to \mathcal{M}$ satisfies
  $G(F(x)) = x$ for all $x$.
\end{theorem}

\begin{proof}
  Non-injectivity gives $x_1 \neq x_2$ with $F(x_1) = F(x_2)$.
  Applying $G$: $G(F(x_1)) = G(F(x_2))$, so exact
  reconstruction of both is impossible.
\end{proof}

\begin{corollary}
  If $E(x) = 0$ for every $x$, then $F$ must be injective.
\end{corollary}

\begin{proposition}[Equivalence-class reconstruction]
  Every reconstruction map acts on equivalence classes,
  not individual states:
  if $y \in [x]$ then $G(F(y)) = G(F(x))$.
\end{proposition}

\begin{proof}
  $y \in [x]$ means $F(y) = F(x)$, so $G(F(y)) = G(F(x))$.
\end{proof}

\begin{definition}[Representational entropy]
  The \emph{representational entropy} of a projection
  $\pi : X \to M$ at $m \in M$ is
  \[
    S_\pi(m) = \log \operatorname{Vol}(\pi^{-1}(m)).
  \]
\end{definition}

\begin{proposition}[Entropy and injectivity]
  \label{prop:entropy-inject}
  $S_\pi(m) = 0$ for all $m$ if and only if $\pi$ is injective.
\end{proposition}

\begin{proof}
  $S_\pi(m) = 0$ iff $\operatorname{Vol}(\pi^{-1}(m)) = 1$
  iff each fibre contains one element iff $\pi$ is injective.
\end{proof}

\begin{theorem}[Minimum reconstruction principle]
  \label{thm:min-recon}
  The reconstruction $G$ minimising expected squared error
  $\mathbb{E}[E(x)^2]$ satisfies
  \[
    G(F(x)) = \frac{\int_{[x]} y\, d\mu(y)}{\mu([x])}.
  \]
  That is, $G$ returns the centroid (Fréchet mean) of the
  equivalence class.
\end{theorem}

\begin{proof}
  Fix $m \in L$.
  Minimise $\int_{F^{-1}(m)} d(y, \hat{y})^2\, d\mu(y)$
  over $\hat{y}$.
  Differentiating and setting to zero yields the Fréchet mean
  condition.
\end{proof}

\begin{theorem}[Projection composition]
  \label{thm:proj-composition}
  For $X \xrightarrow{F} L \xrightarrow{\pi} O$,
  \[
    \ker(F) \subseteq \ker(\pi \circ F).
  \]
\end{theorem}

\begin{proof}
  If $x \in \ker(F)$ then $F(x) = 0$, so
  $(\pi\circ F)(x) = \pi(0)$, hence $x \in \ker(\pi\circ F)$.
\end{proof}

\begin{corollary}
  Every additional observational layer can only increase
  information loss, never decrease it.
\end{corollary}

\begin{definition}[CLIO error functional]
  \[
    \mathcal{E}[F,G]
    = \int_X d(x, G(F(x)))^2\, d\mu(x).
  \]
\end{definition}

\begin{proposition}
  $\mathcal{E}[F,G] \geq 0$, with equality iff $G \circ F = \mathrm{id}$.
\end{proposition}

\begin{theorem}[Existence of optimal reconstruction]
  \label{thm:optimal-recon}
  \statusproven{}
  If $X$ is compact and $d$ is continuous, there exists
  $G^*$ minimising $\mathcal{E}[F,G]$.
\end{theorem}

\begin{proof}
  The space of admissible reconstructions with the topology
  of uniform convergence is compact (Arzelà--Ascoli under
  Lipschitz bounds).
  $\mathcal{E}$ is continuous.
  By the extreme value theorem, a minimiser exists.
\end{proof}


\section{CLIO as constrained optimisation}

\begin{definition}[CLIO problem]
  Given a decomposition $F : \mathcal{M} \to L$ and
  observations $\{m_i\}$, find
  \[
    x^* = \operatorname*{argmin}_{x \in \mathcal{A}}
    \left[\sum_i d(F(x), m_i)^2 + \lambda\,\mathcal{F}[x]\right],
  \]
  where $\mathcal{A}$ is the admissibility class and
  $\mathcal{F}$ is the RSVP free-energy functional.
\end{definition}

The CLIO problem is the variational version of the Minimum
Reconstruction Principle (Theorem~\ref{thm:min-recon}),
extended to include the physical admissibility constraint
$\mathcal{F}[x] < \infty$.

\begin{ChapterConnection}
  The CLIO error functional $\mathcal{E}[F,G]$ introduced
  in the reconstruction propositions of Part~II is the
  objective being minimised here.
  The admissibility constraint $x \in \mathcal{A}$ is the
  sheaf condition from Chapter~\ref{ch:sheaf}.
  CLIO is therefore the intersection of reconstruction theory
  (Part~II) and field theory (Part~IV) mediated by the
  admissibility sheaf (Chapter~\ref{ch:sheaf}).
\end{ChapterConnection}

\begin{ReconView}
  \textbf{Observable:} $m_i = F(x)$ (projected field values).
  \textbf{Hidden:} $x \in \mathcal{A}$ (admissible field state).
  \textbf{CLIO:} Select the admissible $x^*$ that minimises
  mismatch with observations while respecting the energy structure.
  \textbf{Error:} $E(x^*) = d(x^*, G(F(x^*)))$ measures
  how far the optimal reconstruction departs from the truth.
\end{ReconView}

\begin{ProblemSet}
\problem{
  Show that if $F$ is injective and $\mathcal{A}$ contains the
  true state, the CLIO problem has a unique solution with
  $E = 0$.
}
\problem{
  Suppose $\mathcal{A}$ is a closed convex set in $H^k$.
  Show that the CLIO objective $\sum_i d(F(x), m_i)^2
  + \lambda\mathcal{F}[x]$ has at least one minimiser.
}
\problem{
  \textbf{Inverse problem.}
  Given CLIO output $x^*$ and observations $\{m_i\}$,
  compute the residual $\sum_i d(F(x^*), m_i)^2$.
  If this is nonzero, what does it imply about the
  admissibility class $\mathcal{A}$?
}
\end{ProblemSet}
